\section{The largest-eigenvalue estimate}
\label{sec:stw-second-order-proof}

The estimate bounds the complex Hessian in terms of the gradient squared. At a maximum of a function of the largest eigenvalue, the gradient and the potential, the missing torsion indices allow the concavity square to absorb third-derivative terms. Coercivity closes the estimate when the smallest eigenvalue is controlled; in the other case its negative size supplies a positive Hessian square.

Plain subscripts on scalar functions and comma subscripts on matrix coefficients denote coordinate derivatives. The coefficients of the Chern connection are written $\Gamma^q_{ik}$. Choose $\tau>0$ with $Q_0\geq\tau\gamma$.

\Needspace{16\baselineskip}
\begin{theorem}\label{thm:stw-second-order-reconstructed}
Fix smooth Hermitian forms $\gamma_0,\gamma$ on a compact complex manifold of dimension $n\geq2$. For a smooth real function $h$ on $X$, suppose that $v$ is a smooth solution of
\[
 Q=Q_0+P_\gamma(\mathrm i\partial\bar\partial v)+Z(\partial v)>0,
 \qquad \log\det_\gamma Q=h,
\]
with the tensors of \cref{sec:stw-matrix-structure}. If $\|v\|_{C^0}$ and $\|h\|_{C^2}$ are bounded, then
\[
 \sup_X|\mathrm i\partial\bar\partial v|_\gamma
       \leq C\bigl(1+\sup_X|\partial v|_\gamma^2\bigr),
\]
where $C$ depends only on these bounds and the fixed smooth data.
\end{theorem}
\begin{proof}
Adding a constant, arrange $\sup v=0$ and write $-M\leq v\leq0$. Set
\[
 K=1+\sup_X|\partial v|_\gamma^2,\qquad
 \phi(s)=-\tfrac12\log(1-s/(2K)),\qquad
 \psi(t)=D_1\mathrm{e}^{-D_2t}.
\]
The positive constants $D_1,D_2$ will be chosen below using only the fixed data. On $0\leq s\leq K$,
\[
 \frac1{4K}\leq\phi'\leq\frac1{2K},\qquad
 \phi''=2(\phi')^2,\qquad 0\leq\phi\leq\tfrac12\log2.
\]
Let $\lambda_1$ be the largest eigenvalue of $g=\chi+\mathrm i\partial\bar\partial v+W$ relative to $\gamma$. It is positive because $\operatorname{tr}_\gamma g=\operatorname{tr}_\gamma Q>0$. Take a maximum of
\[
 H=\log\lambda_1+\phi(|\partial v|_\gamma^2)+\psi(v).
\]
It suffices to bound $\lambda_1/K$ at this point: the other two terms are uniformly bounded once $D_1,D_2$ are fixed. We may assume that this ratio exceeds any fixed constant specified in the course of the argument.

Choose holomorphic coordinates unitary at the maximum, with $g$ diagonal there and eigenvalues ordered $\lambda_1\geq\cdots\geq\lambda_n$. All estimates below are at this point. Repeated indices are summed only when a sum is displayed. Put $\lambda=\lambda_1$ and use the weights $t_i,f_k$ and their sum $\mathcal F$ from \cref{lem:stw-logdet-coercivity}. Since $\sum_i\lambda_i>0$ and the eigenvalues are ordered,
\[
 |\lambda_i|\leq(n-1)\lambda,\qquad
 |v_{i\bar j}|\leq C\lambda
       \quad\hbox{if }\lambda\geq K.
\]
The constant $C$ may change and is independent of $K$. Until $D_1,D_2$ are fixed below, it is also independent of them. Pure derivatives $v_{ij}$ are not included in this bound.

\paragraph{Differentiating the largest eigenvalue.}
We make its use at a multiple eigenvalue precise. Choose a smooth local unitary frame parallel for the Chern connection at the point, with first vector an eigenvector for $\lambda$. In this frame let $B$ have eigenvalues $0,\epsilon_0,\ldots,\epsilon_0$, where $\epsilon_0>0$ is fixed and small. Then $B\geq0$, $\nabla B=0$ at the point, and its second derivatives are bounded by fixed background constants. The largest eigenvalue $\widetilde\lambda$ of $g-B$ is simple near the point, is at most $\lambda_1(g)$, and equals $\lambda$ there. Hence the corresponding smooth function $\widetilde H$ has a local maximum there. Only the largest eigenvalue needs to be simple.

For a self-adjoint matrix family with a simple top eigenvalue, differentiate its normalized eigenvector equation. Its component in the $p$th eigenspace is the corresponding derivative of the matrix divided by the eigenvalue gap. Differentiating the eigenvalue once more therefore gives
\begin{align*}
 a_k:=\widetilde\lambda_k&=\nabla_k g_{1\bar1},\\
 \widetilde\lambda_{k\bar k}
 &=\nabla_{\bar k}\nabla_k g_{1\bar1}+O(1)
   +\sum_{p>1}\frac{|\nabla_k g_{1\bar p}|^2+
                            |\nabla_k g_{p\bar1}|^2}
                        {\lambda-\lambda_p+\epsilon_0}.
\end{align*}
\Needspace{11\baselineskip}
This formula remains valid if some lower eigenvalues coincide: sum over any orthonormal basis of their eigenspaces. The gaps are positive and at most $n\lambda+\epsilon_0$. In holomorphic coordinates, $a_k=g_{1\bar1,k}+O(\lambda)$, while converting the second covariant derivative to partial derivatives introduces at most
$C\sum_p|g_{1\bar p,k}|+C\lambda$. Convert the off-diagonal covariant derivatives in the positive squares as well; their errors are $O(\lambda)$. Young's inequality absorbs the linear off-diagonal terms. Thus, after decreasing a fixed positive coefficient,
\begin{equation}\label{eq:stw-eigenvalue-basic}
 \widetilde\lambda_{k\bar k}
 \geq g_{1\bar1,k\bar k}
    +\frac1{8n\lambda}\sum_{p>1}
          (|g_{1\bar p,k}|^2+|g_{p\bar1,k}|^2)
    -C(|g_{1\bar1,k}|+\lambda).
\end{equation}
The precise positive constant will not be used.

Write $V_k=\partial_k(|\partial v|_\gamma^2)$. The first derivative of $\widetilde H$ gives
\begin{equation}\label{eq:stw-maximum-first}
 \frac{a_k}{\lambda}=-\phi' V_k-\psi'v_k,
 \qquad
 |V_k|\leq C\sqrt K\sum_p(|v_{pk}|+|v_{p\bar k}|)+CK.
\end{equation}
For the second derivative, the derivatives of the fixed metric in $|\partial v|^2$ cost only a small fraction of its positive Hessian squares and a constant. Therefore
\begin{align}
 \widetilde H_{k\bar k}\geq{}&
  \frac{\widetilde\lambda_{k\bar k}}\lambda
       -\frac{|a_k|^2}{\lambda^2}
       +\phi'\sum_p(v_pv_{\bar p k\bar k}+v_{\bar p}v_{pk\bar k})
       +\frac1{5K}\sum_p(|v_{pk}|^2+|v_{p\bar k}|^2)\notag\\
 &+\phi''|V_k|^2+\psi''|v_k|^2+\psi'v_{k\bar k}-C.
 \label{eq:stw-maximum-second-preliminary}
\end{align}
The third term follows directly by differentiating $\sum\gamma^{p\bar q}v_pv_{\bar q}$ twice and applying $ab\leq\varepsilon a^2+C_\varepsilon b^2$ to each metric-derivative term.

\Needspace{10\baselineskip}
\paragraph{The third derivatives and the torsion term.}
We provide the estimates that remove third derivatives from this inequality. Define
\begin{gather*}
 E=\sum_{k,p} f_k(|v_{pk}|^2+|v_{p\bar k}|^2),\qquad
 S=\sum_{k,p} f_k(|v_{pk}|+|v_{p\bar k}|),\\
 \mathcal C=-D^2\Phi_g(\nabla_1g,(\nabla_1g)^*)\geq0.
\end{gather*}
\Needspace{4\baselineskip}
The first differentiated equation is
\begin{equation}\label{eq:stw-first-differentiated}
 \sum_k f_k\nabla_p g_{k\bar k}=h_p.
\end{equation}
Since $P_\gamma(W)=Z$ and the connection preserves $\gamma$, one has the exact contractions
\begin{equation}\label{eq:stw-contraction-derivatives}
 \sum_k f_k\nabla_p W_{k\bar k}=\sum_i t_i\nabla_p Z_{i\bar i},
 \qquad
 \sum_k f_k\nabla_{\bar1}\nabla_1W_{k\bar k}
       =\sum_i t_i\nabla_{\bar1}\nabla_1Z_{i\bar i}.
\end{equation}
The weights in these identities are evaluated at the point; they are not differentiated.

Here is how the large weights are controlled. The only weight $t_i$ not bounded by $(n-1)f_1$ is $t_1$. In $Z_{1\bar1}$, every gradient index differs from $1$. Thus a pure Hessian in $\nabla_pZ_{1\bar1}$ has an index $q>1$, and can be weighted by $f_q$, which is comparable to $\mathcal F$. Pure derivatives commute in these coordinates. A mixed Hessian can similarly be placed in column $q$ using conjugate symmetry. All other $t_i$ are bounded by a constant times every $f_k$. Derivatives of coefficients cost $C\sqrt K\mathcal F$. Applying these observations to the first identity, and then to \eqref{eq:stw-first-differentiated}, gives, for each $p$,
\begin{equation}\label{eq:stw-weighted-third-bound}
 \left|\sum_k f_kv_{k\bar k p}\right|
       +\left|\sum_k f_kv_{k\bar k\bar p}\right|
       \leq C(S+\sqrt K\mathcal F).
\end{equation}
The connection terms in the differentiated equation contribute weighted mixed Hessians and $C\sqrt K\mathcal F$, already on the right. Bounded derivatives of $h$ are absorbed using the positive lower bound for $\mathcal F$.

For clarity we also track the second derivative of the equation through the eigenvalue estimate. Partial derivatives of $v$ commute, so
\[
 g_{1\bar1,k\bar k}
   =g_{k\bar k,1\bar1}+W_{1\bar1,k\bar k}
          -W_{k\bar k,1\bar1}+O(1).
\]
Twice differentiating the equation covariantly gives
\[
 \sum_k f_k\nabla_{\bar1}\nabla_1g_{k\bar k}
       =h_{1\bar1}+\mathcal C.
\]
The conversion to partial derivatives uses
\[
 \nabla_{\bar1}\nabla_1g_{k\bar k}
       =g_{k\bar k,1\bar1}
           -2\operatorname{Re}\sum_q\overline{\Gamma^q_{1k}}g_{k\bar q,1}
           +O(\lambda).
\]
Rewrite the derivative of $g$ in this term using
$g_{k\bar q,1}=g_{1\bar q,k}+W_{k\bar q,1}-W_{1\bar q,k}+O(1)$ and its conjugate. The first part is absorbed by the off-diagonal squares in \eqref{eq:stw-eigenvalue-basic}, with cost $C\lambda\mathcal F+C\sum_k f_k|g_{1\bar1,k}|$. The $W_{1\bar q,k}$ part is bounded by
$C\lambda\mathcal F+C\sum_{k,p}f_k|v_{pk}|$. For the remaining part use
\[
 2\operatorname{Re}\sum_q\overline{\Gamma^q_{1k}}W_{k\bar q,1}
       =W_{k\bar k,1\bar1}
          -\nabla_{\bar1}\nabla_1W_{k\bar k}+O(\sqrt K).
\]
It cancels the partial second derivative $-W_{k\bar k,1\bar1}$ above. We obtain
\begin{align}
 \sum_k f_k\widetilde\lambda_{k\bar k}
 \geq{}&\mathcal C+\sum_k f_k
       (W_{1\bar1,k\bar k}-\nabla_{\bar1}\nabla_1W_{k\bar k})\notag\\
 &-C\left(\sum_k f_k|g_{1\bar1,k}|
                +\sum_{k,p}f_k|v_{pk}|+\lambda\mathcal F\right).
 \label{eq:stw-eigenvalue-torsion-remainder}
\end{align}

The second contraction in \eqref{eq:stw-contraction-derivatives} bounds its negative torsion term. In the $t_1$ part the possible third derivatives are $v_{q1\bar1}$ and their conjugates with $q>1$; all other parts have weights bounded by $Cf_1$. Commute partial derivatives and use
\[
 v_{1\bar1,k}=g_{1\bar1,k}-\chi_{1\bar1,k}-W_{1\bar1,k}.
\]
The exclusions in \cref{lem:stw-torsion-coefficients} ensure that any remaining pure $v_{11}$ carries only the small weight $Cf_1$. Derivatives of the connection contribute second derivatives with the same weighted bounds. It follows that
\[
 \left|\sum_k f_k\nabla_{\bar1}\nabla_1W_{k\bar k}\right|
 \leq C\left(\sum_k f_k|g_{1\bar1,k}|
             +\sum_{k,p}f_k|v_{pk}|+\lambda\mathcal F\right).
\]
For the positive torsion term write $W_{1\bar1}=\sum_p(c^pv_p+\overline{c^p}v_{\bar p})$ with fixed smooth coefficients. Two derivatives give $c^p\sum_k f_kv_{pk\bar k}$ and its conjugate, plus terms bounded by $C(S+\sqrt K\mathcal F)$. Equation \eqref{eq:stw-weighted-third-bound} bounds the former. In this estimate the mixed Hessians can be bounded by $C\lambda$, giving
\[
 \left|\sum_k f_k W_{1\bar1,k\bar k}\right|
       \leq C\left(\sum_{k,p}f_k|v_{pk}|+\lambda\mathcal F\right).
\]
\Needspace{12\baselineskip}
Substitute these estimates in \eqref{eq:stw-eigenvalue-torsion-remainder}, then in \eqref{eq:stw-maximum-second-preliminary}. The third derivatives in the latter are bounded by \eqref{eq:stw-weighted-third-bound} and $\phi'|\partial v|\leq C/\sqrt K$. The resulting linear Hessian terms are absorbed by a small part of $E/K$, with cost $C\mathcal F$. Thus
\begin{equation}\label{eq:stw-key-maximum-inequality}
 \begin{split}
 0\geq{}&\frac{\mathcal C}{\lambda}
       -\sum_k f_k\frac{|a_k|^2}{\lambda^2}
       +\frac{E}{6K}+\phi''\sum_k f_k|V_k|^2
       +\psi''\sum_k f_k|v_k|^2+\psi'\sum_k f_kv_{k\bar k}\\
 &-C\left(\sum_k f_k\frac{|g_{1\bar1,k}|}{\lambda}+\mathcal F\right).
 \end{split}
\end{equation}

\paragraph{Cancellation when the smallest eigenvalue is controlled.}
Fix $0<\delta<1/2$, to be chosen after $D_1,D_2$. First suppose $\lambda_n\geq-\delta\lambda$. Put $I=\{k:f_k>\delta^{-1}f_1\}$, so $1\notin I$. Equation \eqref{eq:stw-maximum-first} and $\phi''=2(\phi')^2$ give
\begin{align*}
 -\sum_{k\notin I}f_k|a_k|^2/\lambda^2
 &\geq-\phi''\sum_{k\notin I}f_k|V_k|^2
              -2(\psi')^2\delta^{-1}f_1K,\\
 -2\delta\sum_{k\in I}f_k|a_k|^2/\lambda^2
 &\geq-2\delta\phi''\sum_{k\in I}f_k|V_k|^2
              -4\delta(\psi')^2\sum_{k\in I}f_k|v_k|^2.
\end{align*}
\Needspace{7\baselineskip}
Choose $\delta$ small enough that $4\delta(\psi')^2\leq\psi''/2$ throughout $[-M,0]$. By \eqref{eq:stw-concavity-square},
\[
 \frac{\mathcal C}{\lambda}
 \geq\sum_{k\in I}\frac{f_k-f_1}{\lambda(\lambda-\lambda_k)}
                         |\nabla_1g_{k\bar1}|^2
 \geq(1-2\delta)\sum_{k\in I}
                 \frac{f_k|\nabla_1g_{k\bar1}|^2}{\lambda^2}.
\]
Here $(1-\delta)/(1+\delta)\geq1-2\delta$.

\Needspace{7\baselineskip}
We need the following quantitative comparison of the two derivatives. For every $\varepsilon>0$, provided $\lambda/K$ is sufficiently large depending on the fixed function $\psi$,
\begin{equation}\label{eq:stw-derivative-cancellation}
 \sum_{k\in I}f_k\frac{|\nabla_1g_{k\bar1}|^2-|a_k|^2}{\lambda^2}
 \geq-\frac{E}{12K}-C\mathcal F
       -\varepsilon C|\psi'|\mathcal F
       -C_\varepsilon|\psi'|\sum_k f_k|v_k|^2.
\end{equation}
Here is a proof keeping the weights explicit. Let $U=\sum_{p>1,q}|v_{pq}|$. Commuting partial derivatives in $g$ gives, for $k>1$,
\[
 \nabla_1g_{k\bar1}-a_k
       =W_{k\bar1,1}-W_{1\bar1,k}+O(\lambda),\qquad
 |W_{k\bar1,1}|+|W_{1\bar1,k}|\leq C(\lambda+U).
\]
The second inequality uses the absence of $v_1$ in $Z_{k\bar1}$: differentiating in direction $1$ introduces no pure $v_{11}$. All mixed Hessians are bounded by $C\lambda$, and every other pure Hessian belongs to $U$.

\Needspace{10\baselineskip}
Use $|a+e|^2-|a|^2\geq-2|a||e|$, then \eqref{eq:stw-maximum-first}. For $k>1$, the pure Hessians in $V_k$ all belong to $U$. If $H_k=\sum_p|v_{p\bar k}|$, the resulting negative terms, divided by $\lambda^2$, are bounded by a constant times
\[
 K^{-1/2}H_k+K^{-1/2}U
       +\frac{U^2}{\lambda\sqrt K}+1
       +|\psi'||v_k|+\frac{|\psi'|}{\lambda}|v_k|U.
\]
We used $H_k\leq C\lambda$ and $\lambda\geq K$ to simplify the cross terms. The weight comparison \eqref{eq:stw-weight-comparability}, symmetry of pure derivatives and Cauchy--Schwarz give $\mathcal F U^2\leq C E$. Consequently, after multiplication by $f_k$ and summation, the first two terms are bounded by any fixed small fraction of $E/K$ plus $C\mathcal F$. The third is at most $CE/(\lambda\sqrt K)$ and is absorbed by another such fraction for large $\lambda/K$. For the last two terms use
\[
 |v_k|\leq\varepsilon+C_\varepsilon|v_k|^2,\qquad
 |v_k|U\leq |v_k|^2+U^2.
\]
The $U^2$ contribution is at most $C|\psi'|E/\lambda$, again absorbable for large $\lambda/K$. Allocating the fractions to sum to $1/12$ proves \eqref{eq:stw-derivative-cancellation}.

Insert this cancellation and the preceding inequalities into \eqref{eq:stw-key-maximum-inequality}. The remaining derivative error satisfies, by \eqref{eq:stw-maximum-first},
\[
 C\sum_k f_k|g_{1\bar1,k}|/\lambda
 \leq\frac{E}{60K}+C\mathcal F
      +\varepsilon C|\psi'|\mathcal F
      +C_\varepsilon|\psi'|\sum_k f_k|v_k|^2,
\]
after Young's inequality and adjustment of constants. Also the torsion contraction satisfies
\[
 \left|\sum_k f_kW_{k\bar k}\right|
    =\left|\sum_i t_iZ_{i\bar i}\right|
    \leq C\sum_k f_k|v_k|
    \leq\varepsilon C\mathcal F+C_\varepsilon\sum_k f_k|v_k|^2.
\]
\Needspace{9\baselineskip}
The first inequality uses the same missing gradient index in $Z_{1\bar1}$. Since $\psi'<0$ and $v_{k\bar k}=\lambda_k-\chi_{k\bar k}-W_{k\bar k}$, we conclude, decreasing the remaining coefficient of $E$ if necessary, that
\begin{equation}\label{eq:stw-first-case-closure}
 \begin{split}
 0\geq{}& f_1\left(\frac{\lambda^2}{40K}
                      -2(\psi')^2\delta^{-1}K\right)
    +\left(\frac{\psi''}{2}-C_\varepsilon|\psi'|\right)
                           \sum_k f_k|v_k|^2\\
 &-C_0\mathcal F-\varepsilon C_0|\psi'|\mathcal F
       +|\psi'|\sum_k f_k(\chi_{k\bar k}-\lambda_k).
 \end{split}
\end{equation}
In obtaining the first term we used
$|v_{1\bar1}|^2\geq\lambda^2/2-CK$, which follows from $\chi+W=O(\sqrt K)$.

For large $\lambda$, \cref{lem:stw-logdet-coercivity} bounds the last sum below by $\tau\mathcal F/2$. Choose first $\varepsilon$ with $\varepsilon C_0<\tau/4$. Next choose $D_2$ so that $\psi''/2>C_\varepsilon|\psi'|$, and then choose $D_1$ so that $|\psi'|\tau/4>C_0$ on $[-M,0]$. Finally choose $\delta$ as above. All these choices depend only on the fixed data. Equation \eqref{eq:stw-first-case-closure} then forces $\lambda^2\leq80\,\delta^{-1}\sup_{[-M,0]}|\psi'|^2K^2$, proving the desired bound in the first case.

\paragraph{The remaining eigenvalue configuration.}
Suppose now $\lambda_n<-\delta\lambda$, keeping the same constants. Since $f_n\geq\mathcal F/n$ and $|\chi+W|\leq C\sqrt K$,
\[
 \frac{E}{6K}\geq\frac{f_n}{6K}|v_{n\bar n}|^2
       \geq\frac{\delta^2\lambda^2}{12nK}\mathcal F-C\mathcal F.
\]
\Needspace{9\baselineskip}
In \eqref{eq:stw-key-maximum-inequality} discard $\mathcal C$ and the $\psi''$ term. Since $\psi$ is now fixed, $|\psi'\sum_k f_kv_{k\bar k}|\leq C\lambda\mathcal F$. Use $g_{1\bar1,k}=a_k+O(\lambda)$ and Young's inequality to bound its error by
$\frac12\sum_k f_k|a_k|^2/\lambda^2+C\mathcal F$. This leaves
\[
 0\geq-\tfrac32\sum_k f_k|a_k|^2/\lambda^2
       +\frac{\delta^2\lambda^2}{12nK}\mathcal F
       +\phi''\sum_k f_k|V_k|^2-C\lambda\mathcal F.
\]
By \eqref{eq:stw-maximum-first} and the scalar inequality
$\tfrac32|a+b|^2\leq2|a|^2+6|b|^2$,
\[
 \tfrac32\sum_k f_k|a_k|^2/\lambda^2
       \leq\phi''\sum_k f_k|V_k|^2+CK\mathcal F.
\]
As $\lambda\geq K$, cancellation and division by $\lambda\mathcal F>0$ give $\lambda/K\leq C$. This finishes both cases. The maximum property of $H$ propagates the bound to all points, and $|\lambda_i|\leq(n-1)\lambda_1$ then bounds every mixed Hessian entry, as asserted.
\end{proof}
